% ═══════════════════════════════════════════════════════════════════════
\section{Le mode mathématique}
% ═══════════════════════════════════════════════════════════════════════

\begin{frame}[fragile]{Dans le texte, ou hors du texte}
\begin{CodeRendu}[title={Trois façons d'entrer en mode mathématique}]
Le côté $c$ vérifie $a^2 + b^2 = c^2$.

Hors texte, centré :
\[ c = \sqrt{a^2 + b^2} \]

Numéroté :
\begin{equation}
  A = \pi r^2
\end{equation}
\end{CodeRendu}
  \small
  \begin{itemize}
    \item \texttt{\$...\$} : une formule dans la phrase.
    \item \texttt{\textbackslash[ ... \textbackslash]} : une formule importante, sur sa
      propre ligne.
    \item \texttt{equation} : la même chose, avec un numéro qu'on peut citer.
  \end{itemize}
\end{frame}

\begin{frame}[fragile]{L'essentiel (1) : exposants, fractions, racines}
\begin{CodeRendu}[title={Exposants et indices}]
$x^2 \quad x^{10} \quad a_1
 \quad a_{n+1} \quad 2^{-3}$
\end{CodeRendu}
\begin{CodeRendu}[title={Fractions et racines}]
$\frac{3}{4} \quad \dfrac{x+1}{x-2}
 \quad \sqrt{2} \quad \sqrt[3]{27}$
\end{CodeRendu}
  \small Plus d'un caractère en exposant ou en indice : des accolades
  (\texttt{x\^{}10} donne $x^10$, et non $x^{10}$). \cmd{dfrac} : une
  fraction en grand, même dans le texte.
\end{frame}

\begin{frame}[fragile]{L'essentiel (2) : opérations et symboles}
\begin{CodeRendu}[title={Opérations et relations}]
$3 \times 4 \quad 2 \cdot x \quad a \div b$

$x \leq 5 \quad y \geq 0 \quad a \neq b
 \quad \pi \approx 3{,}14$
\end{CodeRendu}
\begin{CodeRendu}[title={Lettres grecques et ensembles}]
$\alpha \quad \beta \quad \theta
 \quad \Delta \quad \pi$

$x \in \mathbb{R} \quad n \in \mathbb{N}$
\end{CodeRendu}
\end{frame}

\begin{frame}[fragile]{L'essentiel (3) : fonctions et texte}
\begin{CodeRendu}[title={Fonctions : droites et bien espacées}]
$\sin(30^\circ) = \frac12
 \quad \log_2 8 = 3 \quad \cos \theta$
\end{CodeRendu}
\begin{CodeRendu}[title={Du texte dans une formule}]
$P = 2(L + l) \text{ où } L
 \text{ est la longueur}$
\end{CodeRendu}
  \small En mode mathématique, les espaces sont ignorés et les lettres
  sont en italique : on écrit \cmd{sin}, jamais \texttt{sin}, et
  \cmd{text\{...\}} pour des mots.
\end{frame}

\begin{frame}[fragile]{Parenthèses et systèmes}
\begin{CodeRendu}[title={Parenthèses qui s'ajustent}]
\[ \left( \frac{a}{b} \right)^2
   = \frac{a^2}{b^2} \]
\end{CodeRendu}
\begin{CodeRendu}[title={Un système d'équations}]
\[ \begin{cases}
     2x + y = 7 \\
     x - y = 2
   \end{cases} \]
\end{CodeRendu}
\end{frame}

\begin{frame}[fragile]{Une démarche alignée}
\begin{CodeRendu}[title={L'environnement align*}]
\begin{align*}
  3x - 5 &= 10 \\
  3x &= 15 \\
  x &= 5
\end{align*}
\end{CodeRendu}
  \small
  \begin{itemize}
    \item \texttt{\&} marque le point d'alignement (ici, avant le signe
      $=$) ; \texttt{\textbackslash\textbackslash} passe à la ligne suivante.
    \item Sans l'étoile, \texttt{align} numérote chaque ligne.
    \item Idéal pour un corrigé : une étape par ligne, toutes alignées.
  \end{itemize}
\end{frame}

\begin{frame}[fragile]{Les nombres à la française}
  \begin{AuTableau}[Piège fréquent]
    \footnotesize En mode mathématique, la virgule est une ponctuation :
    \LaTeX\ ajoute une espace après elle, et $3,14$ s'affiche avec un trou.
    Avec \texttt{siunitx} réglé sur la virgule décimale, \cmd{num} écrit
    les décimaux et groupe les chiffres ; \cmd{qty} ajoute l'unité.
  \end{AuTableau}
\begin{CodeRendu}[title={Trois solutions}]
Mauvais : $3,14$

Accolades : $3{,}14$

Avec siunitx : $\num{3,14}$
et $\num{1234567,89}$

Unités : \qty{12,5}{\centi\metre},
\qty{4}{\centi\metre\squared}
\end{CodeRendu}
\end{frame}

\begin{frame}[fragile]{Trouver un symbole qu'on ne connaît pas}
  \small
  \begin{itemize}
    \item Detexify (\texttt{detexify.kirelabs.org}) : on dessine le
      symbole, le site propose la commande.
    \item L'IA : « Quelle est la commande \LaTeX\ du symbole “environ
      égal” ? »
    \item La documentation d'Overleaf (\texttt{overleaf.com/learn}).
  \end{itemize}
\begin{CodeRendu}[title={Quelques symboles utiles en géométrie}]
$\angle ABC \quad 45^\circ \quad \overline{AB}
 \quad \vec{u} \quad \parallel \quad \perp
 \quad \sim \quad \cong$
\end{CodeRendu}
\end{frame}

\begin{frame}{À vous}
  \begin{Exercice}{Reproduire ces formules}
    \small Dans votre document, écrivez le code qui produit exactement :
    \begin{enumerate}
      \item $f(x) = -2(x - 3)^2 + 1$
      \item L'aire d'un disque de rayon \qty{5}{\centi\metre} est
        $A = \pi r^2 \approx \qty{78,54}{\centi\metre\squared}$.
      \item \( \displaystyle x = \frac{-b \pm \sqrt{b^2 - 4ac}}{2a} \),
        hors texte et centrée.
      \item La démarche $2(x + 4) = 18$, puis $x + 4 = 9$, puis $x = 5$,
        écrite sur trois lignes alignées sur les signes $=$.
      \item $\sin(\theta) = \dfrac{\text{opposé}}{\text{hypoténuse}}$
    \end{enumerate}
    Les solutions sont à la fin de la présentation.
  \end{Exercice}
\end{frame}
